\documentclass{article}
\usepackage[margin=1in]{geometry}
\usepackage{graphicx}

\title{Workflow - EagleNav}
\author{Group 5}
\date{May 2026}

\begin{document}
\maketitle

\section{Overview}
The workflow of EagleNav, a mobile appplication design to improve accessibility and navigation on campus, describes how the application processes
user input, generates routes, provides instructions, and handles rerouting. The process begins when a user selects a destination using the application.
A routing request is then generated and sent to the routing service to provide the optimal path. To improve navigation
reliability, the workflow also includes route deviation detection and rerouting functionality. 

\section{Workflow}

\begin{enumerate}
    \item The user selects a destination through the application interface. The Navigation Controller validates the destination input and sends a routing request to the Routing Controller, which communicates with the Valhalla Service to generate a route.
    \item After the route is successfully generated, the Guidance Controller is activated. The controller returns the route geometry and turn-by-turn navigation instructions needed for active guidance.
    \item The system continuously listens to updates from the Sensor Fusion and Map Matching modules in order to track the user’s current location accurately.
    \begin{itemize}
        \item If the user remains within the acceptable route tolerance, navigation guidance continues normally.
        \item If the user moves outside the allowed deviation threshold, the system prepares to reroute.
    \end{itemize}
    \item When route deviation is detected multiple consecutive times, the Guidance Controller triggers the rerouting process through the Routing Controller. The current implementation uses an approximate 30-meter threshold before generating a new route.
    \item During navigation, the controller continuously updates the user interface with navigation status messages such as:
    \begin{itemize}
        \item ``On Track''
        \item ``Rerouting...''
        \item ``Destination Reached''
    \end{itemize}
    \item Once navigation is completed or cancelled by the user, the Guidance Controller stops the active guidance session and releases navigation resources.
\end{enumerate}

\newpage
\section{Workflow Diagram}

\begin{figure}[h!]
    \centering
    \resizebox{\textwidth}{!}{
    \begin{tabular}{c}
    
    \fbox{User Selects Destination} \\
    $\downarrow$ \\
    \fbox{Navigation Controller Validates Input} \\
    $\downarrow$ \\
    \fbox{Routing Controller \newline Requests Route (Valhalla)} \\
    $\downarrow$ \\
    \fbox{Route Generated} \\
    $\downarrow$ \\
    \fbox{Guidance Controller Activated} \\
    $\downarrow$ \\
    \fbox{Sensor Fusion + Map Matching Updates} \\
    $\downarrow$ \\
    
    \fbox{
    \begin{tabular}{c}
    On Route $\rightarrow$ Continue Guidance \\
    Off Route $\rightarrow$ Trigger Reroute
    \end{tabular}
    } \\
    
    $\downarrow$ \\
    \fbox{UI Updates (On Track / Rerouting / Arrived)} \\
    $\downarrow$ \\
    \fbox{Navigation Ends (Cancel / Arrival)} \\

    \end{tabular}
    }
    \caption{Compact Navigation Workflow}
\end{figure}
\end{document}
